\chapter{Literature Review}\label{ch:litreview}

The present chapter reviews the prior art behind the five technical
strands on which this thesis rests: distributed training,
high-performance GEMM libraries, efficient-attention schemes,
runtimes for serving LLMs, and GPU-kernel profiling. The goals are
three: to record where each strand currently stands, to draw
attention to the gaps the rest of the thesis sets out to close, and
to fix the vocabulary used in the chapters that follow. The order
mirrors the stack itself (Fig.~\ref{fig:intro_stack}): the discussion
begins at the outermost level---the training schedule---and descends
step by step toward kernel-level profiling.

\paragraph{How to read this chapter.} Sections~\ref{sec:lit_dist}
to~\ref{sec:lit_profiling} share a common internal layout. In each
section the classical problem is briefly restated, the principal
systems or papers that speak to it are surveyed, and a pointer is
left to the chapter of the thesis in which that problem is
reopened. The final section, Section~\ref{sec:lit_gap}, consolidates
all of these gaps into one positioning statement, which is then
reused by the contributions list in
Section~\ref{sec:intro_contributions}. A reader already comfortable
with a given sub-area can jump directly to the forward pointer at
the end of the corresponding section.

\section{Distributed Training of Deep Neural Networks}\label{sec:lit_dist}

\subsection{Data Parallelism and Its Refinements}\label{sec:lit_dp}
Of the distributed-training schemes that exist today, data
parallelism (DP) is both the oldest and the conceptually plainest.
Every worker carries a full replica of the model, consumes a distinct
slice of the mini-batch, and at the end of the step joins an
all-reduce that merges the gradients across replicas. The obvious
downsides are two: the model is stored redundantly on every device,
and on every step the entire gradient tensor has to travel across the
network. Ott~et~al.~\cite{ott2018fairseq} demonstrated that with a
carefully designed learning-rate schedule the accuracy of
single-device training can be recovered under large-batch regimes
spread across many GPUs, and You~et~al.~\cite{you2020lars} contributed
the LARS and LAMB optimisers that make extremely large batches
trainable at all.

The engineering realisation of this idea in PyTorch is Distributed
Data Parallel (DDP)~\cite{li2020pytorch}. DDP groups gradient tensors
into communication \emph{buckets} whose size is chosen to match the
granularity at which the backward pass produces them, so that the
reduction of one bucket overlaps with the compute of the next layer.
These reductions ride on a dedicated CUDA stream that plugs directly
into the autograd engine, and they are dispatched around a ring whose
per-worker bandwidth cost is, in the limit, independent of the
number of workers. The well-known $2(P-1)/P$ factor that appears in
the bandwidth expression of a ring all-reduce follows from the
two-phase reduce-scatter/all-gather decomposition; the hop-count
latency term becomes negligible once the parameter tensor is bigger
than a few tens of megabytes, a regime in which every model
considered in this thesis comfortably sits. For these reasons DDP
has settled in as the industry baseline for single-node multi-GPU
training and, accordingly, as the yardstick against which FSDP and
ZeRO are expected to justify the extra collective operations they
introduce.

\subsection{Model Parallelism and Pipeline Parallelism}\label{sec:lit_mp}
When a model can no longer fit on a single accelerator, some form of
model parallelism becomes unavoidable.
Shoeybi~et~al.~\cite{shoeybi2019megatron} put forward Megatron-LM,
a tensor-parallel framework in which the attention and MLP sub-layers
are cut along carefully chosen dimensions so that a whole training
step reduces to a short sequence of all-reduces. Narayanan
et~al.~\cite{narayanan2021efficient} later scaled this design up
towards the trillion-parameter regime by pairing tensor parallelism
with pipeline parallelism on a $3072$-GPU machine, reaching about
$52\%$ of peak throughput.

Pipeline parallelism partitions the model along its depth and places
the resulting stages on different devices. Its characteristic
difficulty is the \emph{pipeline bubble}: idle time that appears
during warm-up at the beginning of a step and again during cool-down
at the end. GPipe~\cite{huang2019gpipe} shrinks this bubble by having
several micro-batches traverse the pipeline simultaneously, at the
cost of storing activations for all of them in flight.
PipeDream~\cite{harlap2018pipedream} interleaves forward and backward
passes---this is the $1\mathrm{F}1\mathrm{B}$ schedule---and
essentially eliminates the bubble at the price of a mild amount of
gradient staleness.

\subsection{Sharded Data Parallelism: FSDP and ZeRO}\label{sec:lit_shard}
In Fully Sharded Data Parallel (FSDP)~\cite{zhao2023pytorch} the
parameters, their gradients, and the optimiser state are all split
across the data-parallel group. Before a given layer enters either
its forward or its backward pass, the parameters of that layer are
all-gathered from their owners, and as soon as the layer finishes
the shards are thrown away, so that the peak per-GPU memory does not
scale with the size of the full model. The design difficulty in
FSDP is the choice of the \emph{unit of sharding}: pick it too coarse
and the all-gathers pile up into bursty traffic; pick it too fine
and the per-call launch overhead begins to dominate. PyTorch FSDP-2
addresses this with the \texttt{FlatParameter} abstraction, which
rolls up every transformer block into a single sharded tensor and in
that way makes the communication granularity coincide with the
computation granularity.

The DeepSpeed ZeRO line~\cite{rajbhandari2019zero, rajbhandari2021zero}
organises the same sharding question into three levels, each of which
reclaims memory from a different component of the training state.
Level one still replicates the parameters on every device but shards
the Adam moment buffers---and, once the FP32 master copy is
included, those moments are where most of the memory actually goes.
Level two adds gradient sharding on top, so that no single worker
ever sees the full gradient tensor; it only ever materialises its
own slice during the reduce-scatter. Level three goes one step
further and shards the parameters themselves, giving each worker
custody of one slice of the model and forcing it to gather the
remaining slices on demand in front of every forward and backward
pass. Level three saturates the information-theoretic lower bound on
per-worker memory; the price is two additional collectives for each
layer---an all-gather before compute and a reduce-scatter after it---
both of which must be hidden behind compute or step-time regresses
badly. ZeRO-Infinity~\cite{rajbhandari2021zero} bolts NVMe offloading
onto this, pushing trillion-parameter training down to $512$-GPU
boxes. DeepSpeed itself~\cite{rasley2020deepspeed,
rajbhandari2022deepspeed} bundles ZeRO with gradient compression,
loss scaling, and Mixture-of-Experts support, and has become the
default system in many production LLM training pipelines.

Above roughly $10$B parameters, Korthikanti~et~al.\
\cite{korthikanti2022activation} observed that it is no longer
parameter memory but \emph{activation} memory that becomes binding,
and they introduced selective activation recomputation---a
refinement of an older recipe from
Chen~et~al.~\cite{chen2016training} applied to far smaller models.
In the published literature, comparisons of DP, DDP, FSDP, and ZeRO
tend to appear one system at a time rather than as a single
controlled single-node study; Chapter~\ref{ch:distributed} fills that
gap by pinning the model, the dataset, and the hardware and
comparing the four strategies directly.

\subsection{Mixed Precision and Optimiser
Considerations}\label{sec:lit_mp_amp}
Micikevicius~et~al.~\cite{micikevicius2018amp} established that
training weights in FP16 while keeping an FP32 master copy, with
loss scaling adjusted dynamically, preserves convergence behaviour
but roughly halves the memory demand and doubles the Tensor-Core
throughput. BF16 on A100 and FP8 on H100~\cite{nvidia_hopper} extend
this idea yet further; each step introduces a small numerical gap
that has to be contained at the individual layer, usually by
deciding carefully where the accumulation is carried out in higher
precision.

\subsection{Collective Communication Libraries}\label{sec:lit_nccl}
Performance in distributed training ultimately rests on the
collective-communication layer beneath the framework. NCCL~\cite{nccl}
supplies both ring and tree all-reduces on top of NVLink, NVSwitch,
and InfiniBand, and automatically selects among them based on
message size. Castell{\'o}~et~al.~\cite{castello2021accelerating}
investigate a pipelined all-reduce built on MPI, and report that
overlapping the reduce and broadcast phases with intra-node copies
can reach up to a $30\%$ end-to-end gain on commodity clusters.

\section{High-Performance GPU Kernels and GEMM Libraries}\label{sec:lit_gemm}

Every layer of a transformer ultimately stands or falls on the
throughput of a matrix multiply. Vendor libraries such as cuBLAS and
cuBLASLt ship carefully tuned dense GEMM kernels, but their coverage
of the non-canonical shapes, fused epilogues, and grouped variants
that a modern LLM actually demands is patchy. CUTLASS~\cite{cutlass},
by contrast, is a C++ template library that exposes GEMM as a
configurable stack of building blocks, with the tile sizes, the
software pipelining depth, and the Tensor-Core instructions all
placed under explicit user control. Its CuTe sub-library generalises
the underlying layout algebra so that a single template can be
instantiated against many hardware generations and data types.

\subsection{From cuBLAS to CUTLASS}\label{sec:lit_cublas_cutlass}
Because cuBLAS is a closed-source library tuned by NVIDIA around a
small number of canonical GEMM shapes and precisions, it often falls
back to a generic code path on the very shapes that LLMs produce in
practice---skinny, batched, and grouped GEMMs at the heart of
attention and MoE layers---leaving the Tensor Cores noticeably
under-utilised. CUTLASS instead exposes the three-level tile
hierarchy (threadblock, warp, and instruction) directly, so that a
kernel author can align the template with the shape they really
have; the same stack composes with epilogue templates that let bias,
activation, or residual additions be folded into the writeback of
the main GEMM.

\subsection{Triton and the DSL
Approach}\label{sec:lit_triton_dsl}
Triton~\cite{tillet2019triton} follows a rather different route. It
is a Python-embedded domain-specific language that puts the tile
abstraction squarely in the user's hands, in exchange for a modest
fraction of peak throughput relative to a hand-tuned CUDA kernel,
and for dramatically shorter iteration cycles. Much of the later
work in FlashAttention-2~\cite{dao2023flashattention2} and
FlashAttention-3~\cite{dao2024flashattention3} sits on top of these
Triton primitives, and a large part of the serving stack discussed
in the latter half of this thesis is also written in Triton.
Direct, side-by-side comparisons of the Triton and CUTLASS
ecosystems on the exact kernels produced by a transformer forward
pass are still uncommon in the literature; that gap is one of the
things Chapters~\ref{ch:cutlass} and~\ref{ch:triton} are designed
to close.

\subsection{Hardware-Specific Primitives}\label{sec:lit_hw}
The A100 whitepaper~\cite{nvidia_a100} sets out the third-generation
Tensor Cores, rated at $312$\,TFLOPS of FP16, together with the new
\texttt{mma.m16n8k16} and \texttt{mma.m16n8k8} instructions and the
asynchronous \texttt{cp.async} copy path. The Hopper
whitepaper~\cite{nvidia_hopper} extends this hardware vocabulary with
the Tensor Memory Accelerator (TMA), fourth-generation Tensor Cores,
a dedicated transformer engine that natively supports FP8, and the
DSMEM mechanism that lets threadblock clusters cooperate through
shared memory. Behind a layout-algebra interface, CUTLASS~3.x hides
these architectural differences, so that one kernel template can be
retargeted to different architectures by swapping the appropriate
\texttt{Sm}-specific instantiation.

\section{Efficient Attention Mechanisms}\label{sec:lit_attn}

\subsection{Exact Attention Acceleration}\label{sec:lit_flashattn}
The key observation behind the Flash Attention family is that the
$N \times N$ score matrix is ultimately an intermediate quantity and
does not, at any point, have to live in HBM. Provided that queries
are streamed against tiles of keys and values, and that a running
row maximum and row sum are carried from tile to tile through a
correction factor, the final output can be reconstructed without
ever materialising the dense scores in global memory. Using exactly
this trick, the original Flash Attention
kernel~\cite{dao2022flashattention} brings the memory footprint down
from $\mathcal{O}(N^2)$ to $\mathcal{O}(N)$, while still being
bit-identical to the dense reference up to the order of
floating-point additions. FlashAttention-2
\cite{dao2023flashattention2} rebalances how the work is spread
along the sequence axis, so that A100 occupancy no longer collapses
at long context lengths, and
FlashAttention-3~\cite{dao2024flashattention3} rewrites the same
algorithm on top of Hopper's WGMMA and FP8 pipelines, exploiting the
asynchrony between the TMA copy engine and the Tensor-Core
pipeline. Across this thesis, the three kernels serve as the
reference point: any new attention kernel has to match them on
correctness and to beat them either on wall-clock speed or on
sparsity.

\subsection{Sparse and Local Attention}\label{sec:lit_sparse}
A different line of work tackles the quadratic cost at the level of
the definition of attention itself, by throwing away the
dependencies between tokens that are sufficiently far apart. The
Sparse Transformer~\cite{child2019generating} was the first design
in this family to be widely adopted; it stitches a strided pattern
together with a fixed pattern to handle sequences in the tens of
thousands of tokens while keeping the per-layer cost linear in $N$.
Longformer~\cite{beltagy2020longformer} simplifies the design to a
local sliding window supplemented by a small set of global tokens
for task-specific anchoring, and BigBird~\cite{zaheer2020bigbird}
restores the universal-approximation property by layering a random
long-range term on top of the local window. The open-weight models
of the current generation have converged on the simplest of these
patterns---a plain sliding window a few thousand positions wide---
because empirically the softmax weight outside that radius is
already dominated by numerical noise. It is precisely this
architectural convergence that turns a window-aware kernel from a
curiosity into something of practical commercial value.

A complementary class of methods replaces full attention with a
cheaper-to-compute surrogate. Performer~\cite{choromanski2020performer}
and Linformer~\cite{wang2020linformer} linearise the attention kernel;
Reformer~\cite{kitaev2020reformer} leans on low-rank hashing. Each of
these trades asymptotic complexity for a small accuracy gap, and in
practice they have seen only moderate adoption in production
settings. The sliding-window path taken in Chapter~\ref{ch:swfa}
keeps the attention numerics exact and exploits the sparsity at the
kernel level, rather than rewriting the mathematical definition.

\subsection{Kernel-Level Fusions for Attention}\label{sec:lit_fused}
Another family of attention variants is tailored specifically to
inference. PagedAttention~\cite{vllm} keeps the key-value cache as a
pool of fixed-size pages, so that the attention kernel accesses the
cache through 32-bit page indices and fragmentation is eliminated.
RadixAttention~\cite{sglang} generalises that idea to a radix tree,
so that prompt prefixes shared between concurrent requests are kept
only once. Multi-Query Attention (MQA)~\cite{shazeer2019mqa} and
Grouped-Query Attention (GQA)~\cite{ainslie2023gqa} shrink the
KV cache by a factor of $H$ (the total head count) or $H/G$ (the
group count), which in turn cuts the memory bandwidth demanded by
the decode step. Each of these designs changes the memory layout
the attention kernel has to see, and therefore has direct
consequences for kernel scheduling; those consequences are taken
up in Chapter~\ref{ch:sglang}.

\subsection{Positional Encodings for Long
Context}\label{sec:lit_rope}
Long-context attention cannot really be discussed apart from the
positional encoding it carries. Rotary position embeddings
(RoPE)~\cite{su2021rope} extrapolate gracefully to sequences longer
than those seen during training, and when combined with dynamic
NTK-aware scaling they allow a model originally trained at
$N = 4{,}096$ to be served at $N = 32{,}768$ with only a mild
degradation. SWFA composes directly with RoPE and needs no
modification, because the window constraint operates entirely on
the attention score matrix.

\subsection{Block Sparsity as a Design
Pattern}\label{sec:lit_block_sparsity}
A theme running through all of the above is that \emph{block-level}
sparsity is practical while \emph{element-level} sparsity is not.
The reason is structural: modern Tensor Cores issue the same
$16 \times 8 \times 16$ MMA instruction irrespective of how many
elements of the operand tile happen to be zero, so element-level
sparsity produces no wall-clock saving unless the hardware itself
supports a structured pruning format. Block sparsity, by contrast,
aligns with the threadblock-level tiling of CUTLASS and CuTe in the
natural way: a block labelled ``skip'' is simply never entered by
the loop nest of Algorithm~\ref{alg:swfa}. That observation is
exactly why SWFA is able to extract measurable speedups without a
custom Tensor-Core ISA. The same idea has been applied to the
strided sparsity of the Sparse
Transformer~\cite{child2019generating} and to the global-plus-random
mask of BigBird~\cite{zaheer2020bigbird}; to the best of our
knowledge, however, it has not yet been combined with Flash
Attention's online softmax inside a single published kernel.

\section{LLM Serving Systems and KV-Cache Management}\label{sec:lit_serving}

The workload profile of serving differs qualitatively from that of
training. Requests come in asynchronously, prompts share large
prefixes between one another, and---unlike training---the decode
step is bound by memory bandwidth rather than by compute. A single
chat request, for instance, may add only a few hundred tokens of
unique suffix on top of a system prompt that is several thousand
tokens long and is shared with every other request active at the
same time; any serving runtime that does not exploit that shared
prefix pays a cost linear in the size of the prefix, for every
request. vLLM~\cite{vllm} was the first runtime to see wide adoption
that treated the KV cache as a pool of fixed-size pages and
introduced \emph{continuous batching}, in which new requests can
join the running batch at any decode step. Building on this,
SGLang~\cite{sglang} contributes both RadixAttention and a
structured-generation language that compiles complex prompt flows
into a scheduled dataflow of kernel calls. Both systems stress the
kernel layer in novel ways: RadixAttention alters the access pattern
that the attention kernel has to respect, and continuous batching
produces a stream of many small, tightly scheduled kernel launches.

The kernel-level contributions of Chapter~\ref{ch:sglang} sit within
this lineage. Where the vLLM and SGLang papers describe their work
at the system level, this thesis zooms in on the \emph{individual
kernels} that the runtime dispatches to, and reports their
performance in isolation so that every improvement can be ascribed
to a specific kernel change.

\section{GPU Kernel Profiling and Auto-tuning}\label{sec:lit_profiling}

Two traditions have shaped kernel-level performance analysis. On
one side, vendor metric collectors---Nsight
Compute~\cite{nsight} on current NVIDIA silicon, and \texttt{nvprof}
before it---surface occupancy, cache hit rates, and the
instruction-mix breakdowns that are needed to pin a stall to a
particular hardware unit. On the other side, the roofline
model~\cite{williams2009roofline} provides an analytical upper
bound: a kernel's sustained throughput is limited by whichever
comes first, the peak compute rate of the device or the product of
the kernel's arithmetic intensity with the memory bandwidth. On an
A100 running in FP16---$312$\,TFLOPS of Tensor-Core throughput
against $1{,}555$\,GB/s of HBM2e bandwidth---the two ceilings
intersect around an arithmetic intensity of roughly $200$\,FLOP per
byte. Dense square GEMMs comfortably clear that line and therefore
sit on the compute ceiling; attention-score GEMMs with a small
$d_{\text{head}}$, and the single-query decode kernels that show up
in serving, sit well below it and are bandwidth-bound by
construction.

Triton exposes a built-in \texttt{autotune} mechanism that sweeps
over block shapes, warp counts, and pipeline depths and caches the
winning configuration keyed by input shape. CUTLASS offers an
analogous \texttt{GemmUniversal} enumerator along with its
\texttt{cutlass\_profiler} binary. The open literature is rich in
single-kernel auto-tuning studies but comparatively thin in
end-to-end pipelines that cross-validate Triton, CUTLASS, and
Torch-inductor on the same set of LLM kernels---an
end-to-end pipeline of exactly that kind is presented in
Chapter~\ref{ch:triton}.

\section{Energy and Carbon Considerations}\label{sec:lit_energy}
According to
Patterson~et~al.~\cite{patterson2021carbon}, the training of GPT-3
released on the order of $552$\,t\,CO\textsubscript{2}e---
comparable to $115$ transatlantic passenger flights---and the
cumulative cost of serving a model over its deployment lifetime can
end up exceeding the one-off cost of training it. That observation
motivates two choices in this thesis: the comparison of distributed
training strategies in Chapter~\ref{ch:distributed} is energy-aware,
and every result table in that chapter reports kWh per epoch next
to the familiar throughput and memory metrics.

\section{Cross-Cutting Concerns}\label{sec:lit_cross}

Before the positioning section it is worth discussing three issues
that cut across every one of the five layers above.

\subsection{Numerical Precision at
Scale}\label{sec:lit_precision}
Mixed-precision training (Section~\ref{sec:lit_mp_amp}) is only
reliable when every kernel in the stack agrees on the accumulator
precision. Both CUTLASS and Triton let the accumulator type be
chosen independently of the input data type, and the
distributed-training frameworks---DDP, FSDP, and DeepSpeed---expose
a \texttt{reduce\_dtype} that keeps the cross-device reduction in
FP32 even when the parameters themselves are in FP16. This thesis
sticks to the same convention at every layer: FP16 or BF16 inputs,
FP32 accumulators, and FP32 reductions. The numerical validation
performed in Chapter~\ref{ch:swfa} confirms that this choice is
enough to match a vanilla PyTorch reference to within $10^{-4}$
tolerance in FP32 and $10^{-2}$ in FP16. Looking ahead, FP8 on
Hopper~\cite{nvidia_hopper} and the emerging FP4/NVFP6 formats will
shift the accumulator-precision question out of the main MMA loop
and into the epilogue; the CUTLASS epilogue templates built in
Chapter~\ref{ch:cutlass} are structurally prepared for this, though
the shape-indexed auto-tune cache of Chapter~\ref{ch:triton} will
need to grow an additional accumulator-precision axis.

\subsection{Reproducibility in a Non-Deterministic
Stack}\label{sec:lit_repro}
By default, every layer of the stack is non-deterministic.
PyTorch's own reduction kernels choose their reduction order based
on the launch configuration, NCCL chooses between ring and tree
algorithms based on the message size, and SGLang's continuous
batching produces an ordering of kernel launches that depends on
which requests happen to be live at the moment. Strict bit-level
reproduction of a run is therefore out of reach. The convention
adopted here is the weaker but practical one of \emph{statistical}
reproducibility: the random seeds and the hardware are fixed, and
every reported configuration carries a relative standard deviation
below $1\%$. Wherever a deterministic path does exist---in
particular for the SWFA numerical-accuracy validation of
Chapter~\ref{ch:swfa}---it is taken, and the deterministic output
is compared bit-for-bit against the PyTorch reference.

\subsection{Energy as a First-Class
Metric}\label{sec:lit_energy_metric}
In the spirit of
Patterson~et~al.~\cite{patterson2021carbon}, energy per epoch is
reported in this thesis alongside throughput, at every layer where
it can be measured. The reason is not, in the first instance, an
environmental one. At data-centre scale, energy is simply a proxy
for the sum of idle time, memory stalls, and under-utilised
compute---costs that a raw throughput number does not expose. A
training run that completes in $8$ hours but consumes $20\%$ more
energy than strictly necessary has $20\%$ of wasted work parked
somewhere in the stack, and identifying which layer is hosting that
waste is exactly what the controlled comparison in
Chapter~\ref{ch:distributed}, paired with the cross-validating
profiling pipeline of Chapter~\ref{ch:triton}, is designed to do.

\section{Positioning of This Thesis}\label{sec:lit_gap}

The impression left by the literature is one of rapid, localised
progress in each of the five areas covered above---distributed
training, GEMM libraries, efficient attention, serving runtimes,
and kernel profiling---together with rather limited integration
across them. Distributed-training studies rarely report
kernel-level measurements; papers on attention kernels are
typically benchmarked in isolation; and system-level papers on
serving frameworks tend to treat the attention kernel as a fixed
dependency. The view taken here is the complementary one: the
training, kernel, attention, serving, and profiling layers have to
be co-designed rather than optimised in isolation.
\begin{itemize}
    \item A sliding-window attention kernel only pays for itself if
          the serving runtime actually exposes the window to it, and
          if the GEMM library underneath can launch the resulting
          batched and sparse tiles efficiently.
    \item A distributed-training schedule only reaches its peak
          efficiency when the GEMM and attention kernels beneath
          it are themselves well tuned; any Tensor Cores that are
          under-utilised at the kernel level are inherited as idle
          Tensor Cores at every scale of parallelism.
    \item The only way to reason quantitatively about these
          cross-layer interactions, without being misled by
          per-library calibration differences, is to use a
          profiling pipeline that reports a common set of metrics
          across Triton, CUTLASS, and PyTorch.
\end{itemize}
Chapters~\ref{ch:distributed}--\ref{ch:triton} take these layers
up one by one, and the evaluation at the end of each chapter
demonstrates how they compose across the stack.

Table~\ref{tab:lit_summary} collects the principal systems cited in
the chapter, recording for each system the layer of the stack at
which it intervenes and the primary axis along which it contributes.

\begin{table}[h]
\centering
\caption{Principal prior systems organised by stack layer.}
\label{tab:lit_summary}
\begin{tabular}{lllc}
\toprule
\textbf{System} & \textbf{Layer} & \textbf{Primary axis} & \textbf{Reference} \\
\midrule
DDP              & Distributed & Bucketed all-reduce  & \cite{li2020pytorch} \\
Megatron-LM      & Distributed & Tensor parallelism    & \cite{shoeybi2019megatron} \\
GPipe            & Distributed & Pipeline parallelism  & \cite{huang2019gpipe} \\
PipeDream        & Distributed & 1F1B pipeline         & \cite{harlap2018pipedream} \\
FSDP             & Distributed & Parameter sharding    & \cite{zhao2023pytorch} \\
ZeRO             & Distributed & Three-stage sharding  & \cite{rajbhandari2019zero} \\
ZeRO-Infinity    & Distributed & NVMe offload          & \cite{rajbhandari2021zero} \\
DeepSpeed        & Distributed & System integration    & \cite{rasley2020deepspeed} \\
CUTLASS          & Kernel      & Template GEMM         & \cite{cutlass} \\
Triton           & Kernel      & Tile-level DSL        & \cite{tillet2019triton} \\
Flash Attention  & Attention   & Online softmax        & \cite{dao2022flashattention} \\
FlashAttention-2 & Attention   & Work partitioning     & \cite{dao2023flashattention2} \\
FlashAttention-3 & Attention   & WGMMA + FP8           & \cite{dao2024flashattention3} \\
Longformer       & Attention   & Local + global        & \cite{beltagy2020longformer} \\
BigBird          & Attention   & Random + local        & \cite{zaheer2020bigbird} \\
MQA / GQA        & Attention   & Head reduction        & \cite{shazeer2019mqa,ainslie2023gqa} \\
vLLM             & Serving     & Paged KV cache        & \cite{vllm} \\
SGLang           & Serving     & Radix KV + DSL        & \cite{sglang} \\
Nsight Compute   & Profiling   & Metric collection     & \cite{nsight} \\
Roofline         & Profiling   & Analytical model      & \cite{williams2009roofline} \\
\bottomrule
\end{tabular}
\end{table}
